% GENERATED FILE. Edit canonical lessons, then run npm run generate:books.
% canonical-source-hash: fnv1a64:a7a6c67dc105a370
% canonical-lessons: BN-C242-berono, BN-C242-chhorano, BN-C242-gotano, BN-C242-mora2, BN-C242-khora

\chapter{To go out, To scatter, To roll up, To wrap, To dig}
\label{ch:bn-to-go-out-to-scatter-to-roll-up-to-wrap-to-dig}
\hlchaptermodality{242}

\begin{chapteropening}
\textbf{By the end of this chapter:} \emph{I can say to go out, to scatter, to roll up, to wrap and to dig.}

Complete the last lesson of chapter 242: I can say to go out, to scatter, to roll up, to wrap and to dig.
\end{chapteropening}

\section[\texorpdfstring{\bn{বেরোনো} (berono)}{বেরোনো (berono)}]{\bn{বেরোনো} (berono) — to go out}

\label{lesson:BN-C242-berono}



\begin{quote}
Before the new one: say the Bengali for to inform, then the Bengali for to chat.
\end{quote}

\subsection*{You'll want to know: \bn{বেরোনো}}
\textbf{\bn{বেরোনো}} — \emph{berono} — \textquotedblleft{}to go out\textquotedblright{}.

\begin{cousinweb}
The origin of \textbf{\bn{বেরোনো}} (\emph{berono}) is not settled: CDIAL has no entry for it, and a search by form and by the glosses \textquotedblleft{}go out\textquotedblright{} and \textquotedblleft{}come out\textquotedblright{} finds only words of other shapes, such as Bengali \emph{nikalā}, \textquotedblleft{}to come out\textquotedblright{} (CDIAL 7478).
\end{cousinweb}

\begin{grammarlens}[title={What you've built}]
One more word for this chapter.
\end{grammarlens}

\subsection*{Your turn}
\begin{itemize}
\raggedright
  \item \emph{Say it:} \emph{berono}
  \item \emph{Say it:} \emph{berono}, once more
  \item \emph{Say it:} say \emph{berono}
  \item \emph{Recall:} say the Bengali for to inform, then the Bengali for to chat, then say \emph{berono} again
\end{itemize}

\begin{tcolorbox}[breakable,colback=teal!4,colframe=teal!35!black,title={Before you move on}]
What does \bn{বেরোনো} mean? (\textquotedblleft{}To go out\textquotedblright{}.) Say it once more.
\end{tcolorbox}
\section[\texorpdfstring{\bn{ছড়ানো} (chhôṛāno)}{ছড়ানো (chhôṛāno)}]{\bn{ছড়ানো} (chhôṛāno) — to scatter, to spread}

\label{lesson:BN-C242-chhorano}



\begin{quote}
Before the new one: say the Bengali for to chat, then the Bengali for to go out.
\end{quote}

\subsection*{You'll want to know: \bn{ছড়ানো}}
\textbf{\bn{ছড়ানো}} — \emph{chhôṛāno} — \textquotedblleft{}to scatter, to spread\textquotedblright{}.

\begin{cousinweb}
\textbf{\bn{ছড়ানো}} (\emph{chhôṛāno}) is a tadbhava, a word inherited and reshaped by sound change. It comes from \emph{*chaṭ}, \textquotedblleft{}scatter, sift\textquotedblright{}, a form CDIAL reconstructs from its descendants (CDIAL 4965). CDIAL lists Bengali \emph{chaṛāna}, \textquotedblleft{}to scatter\textquotedblright{}.
\end{cousinweb}

\begin{grammarlens}[title={What you've built}]
One more word for this chapter.
\end{grammarlens}

\subsection*{Your turn}
\begin{itemize}
\raggedright
  \item \emph{Say it:} \emph{chhôṛāno}
  \item \emph{Say it:} \emph{chhôṛāno}, once more
  \item \emph{Say it:} say \emph{chhôṛāno}
  \item \emph{Recall:} say the Bengali for to chat, then the Bengali for to go out, then say \emph{chhôṛāno} again
\end{itemize}

\begin{tcolorbox}[breakable,colback=teal!4,colframe=teal!35!black,title={Before you move on}]
What does \bn{ছড়ানো} mean? (\textquotedblleft{}To scatter, to spread\textquotedblright{}.) Say it once more.
\end{tcolorbox}
\section[\texorpdfstring{\bn{গোটানো} (goṭāno)}{গোটানো (goṭāno)}]{\bn{গোটানো} (goṭāno) — to roll up}

\label{lesson:BN-C242-gotano}



\begin{quote}
Before the new one: say the Bengali for to go out, then the Bengali for to scatter.
\end{quote}

\subsection*{You'll want to know: \bn{গোটানো}}
\textbf{\bn{গোটানো}} — \emph{goṭāno} — \textquotedblleft{}to roll up\textquotedblright{}.

\begin{cousinweb}
\textbf{\bn{গোটানো}} (\emph{goṭāno}) appears to be built on Bengali \emph{goṭā}, \textquotedblleft{}whole, a seed\textquotedblright{}, which CDIAL lists under \emph{*gōṭṭa}, \textquotedblleft{}something round\textquotedblright{} (CDIAL 4271): \textquotedblleft{}to make round\textquotedblright{}.
\end{cousinweb}

\begin{grammarlens}[title={What you've built}]
One more word for this chapter.
\end{grammarlens}

\subsection*{Your turn}
\begin{itemize}
\raggedright
  \item \emph{Say it:} \emph{goṭāno}
  \item \emph{Say it:} \emph{goṭāno}, once more
  \item \emph{Say it:} say \emph{goṭāno}
  \item \emph{Recall:} say the Bengali for to go out, then the Bengali for to scatter, then say \emph{goṭāno} again
\end{itemize}

\begin{tcolorbox}[breakable,colback=teal!4,colframe=teal!35!black,title={Before you move on}]
What does \bn{গোটানো} mean? (\textquotedblleft{}To roll up\textquotedblright{}.) Say it once more.
\end{tcolorbox}
\section[\texorpdfstring{\bn{মোড়া} (moṛā)}{মোড়া (moṛā)}]{\bn{মোড়া} (moṛā) — to wrap}

\label{lesson:BN-C242-mora2}



\begin{quote}
Before the new one: say the Bengali for to scatter, then the Bengali for to roll up.
\end{quote}

\subsection*{You'll want to know: \bn{মোড়া}}
\textbf{\bn{মোড়া}} — \emph{moṛā} — \textquotedblleft{}to wrap\textquotedblright{}.

\begin{cousinweb}
\textbf{\bn{মোড়া}} (\emph{moṛā}), \textquotedblleft{}to wrap\textquotedblright{}, appears to be the verb CDIAL lists as Bengali \emph{moṛā}, \textquotedblleft{}to twist\textquotedblright{}, under Sanskrit \emph{mōṭati}, a form of \emph{muṭati}, \textquotedblleft{}twists\textquotedblright{} (CDIAL 10186).
\end{cousinweb}

\begin{grammarlens}[title={What you've built}]
One more word for this chapter.
\end{grammarlens}

\subsection*{Your turn}
\begin{itemize}
\raggedright
  \item \emph{Say it:} \emph{moṛā}
  \item \emph{Say it:} \emph{moṛā}, once more
  \item \emph{Say it:} say \emph{moṛā}
  \item \emph{Recall:} say the Bengali for to scatter, then the Bengali for to roll up, then say \emph{moṛā} again
\end{itemize}

\begin{tcolorbox}[breakable,colback=teal!4,colframe=teal!35!black,title={Before you move on}]
What does \bn{মোড়া} mean? (\textquotedblleft{}To wrap\textquotedblright{}.) Say it once more.
\end{tcolorbox}
\section[\texorpdfstring{\bn{খোঁড়া} (khõṛā)}{খোঁড়া (khõṛā)}]{\bn{খোঁড়া} (khõṛā) — to dig}

\label{lesson:BN-C242-khora}



\begin{quote}
Before the new one: say the Bengali for to roll up, then the Bengali for to wrap.
\end{quote}

\subsection*{You'll want to know: \bn{খোঁড়া}}
\textbf{\bn{খোঁড়া}} — \emph{khõṛā} — \textquotedblleft{}to dig\textquotedblright{}.

\begin{cousinweb}
\textbf{\bn{খোঁড়া}} (\emph{khõṛā}) is a tadbhava, a word inherited and reshaped by sound change. It comes from \emph{*khodd}, \textquotedblleft{}dig\textquotedblright{}, a form CDIAL reconstructs from its descendants (CDIAL 3934). CDIAL lists Bengali \emph{khõṛā}, \textquotedblleft{}to dig\textquotedblright{}, here and also under \emph{*khuṭati}, \textquotedblleft{}breaks\textquotedblright{} (CDIAL 3892).
\end{cousinweb}

\begin{grammarlens}[title={What you've built}]
One more word for this chapter.
\end{grammarlens}

\subsection*{Your turn}
\begin{itemize}
\raggedright
  \item \emph{Say it:} \emph{khõṛā}
  \item \emph{Say it:} \emph{khõṛā}, once more
  \item \emph{Say it:} say \emph{khõṛā}
  \item \emph{Recall:} say the Bengali for to roll up, then the Bengali for to wrap, then say \emph{khõṛā} again
\end{itemize}

\begin{tcolorbox}[breakable,colback=teal!4,colframe=teal!35!black,title={Before you move on}]
What does \bn{খোঁড়া} mean? (\textquotedblleft{}To dig\textquotedblright{}.) Say it once more.
\end{tcolorbox}
